\documentclass[12pt]{article}
%\renewcommand\normalsize{\fontsize{10pt}{12pt}\selectfont}
\usepackage{graphicx} % Required for inserting images
\usepackage[left=2cm,right=2cm,top=2cm,bottom=1.5cm]{geometry}
\usepackage{wrapfig}
%\pagestyle{empty} 

\title{Quiz 4}
%\author{Dennis Wilkinson}
%\date{March 31 2025; due Monday April 7}

\begin{document}


\setlength{\parindent}{0pt}
\setlength{\parskip}{3mm}

FS1 Wilkinson Quiz 4, May 22, 2025

\textbf{1.} Consider the following procedure, as illustrated below: (1) a positively charged insulator is brought near a neutral metal sphere A; (2) a second neutral sphere B is touched to the first; (3) the spheres are separated; (4) the insulator is removed.

\vspace{1in}

\textbf{a.} Is there charge on the sphere(s)? If so, what sign on each, and why? If not, why not?
\vspace{1.5in}

\textbf{b.} When charge moved around in this process, what moved - circle one: 
(a) electrons
(b) protons
(c) ions
(d) some combination of these.

\bigskip

\textbf{2.} Two metal objects are fixed in place and have charge +Q and -Q, as shown.
\vspace{1.75in}

\textbf{a.} Write in the \textit{direction} of the electric field at points A, B, and C (draw arrows on the picture at these points)

\textbf{b.} If the magnitude of the field at point A is 20000 N/C, and a $-30$ nC (that is, $-30 \times 10^{-9}$ C) charge is placed there, what force - magnitude and direction - does it feel?

\vspace{1 in}

\textbf{c.} Is the 20000 N/C electric field at A enough to ionize the air and cause a shock if, say, you put your finger there?

\end{document}